%===============================================================================
% APPENDIX XIX: LIT-LEARNING
% Experience Goods, Belief Updating, and Learning Literature
%===============================================================================
% Category: LIT
% Version: 1.0
% Date: January 2026
% Language: english
% Dependencies: PREDICT-HABIT (HF), FORMAL-PROBABILITY (PBB), CORE-AWARE (AU)
%===============================================================================

\documentclass[11pt, a4paper]{article}
\usepackage[utf8]{inputenc}
\usepackage[english]{babel}
\usepackage[margin=1in]{geometry}
\usepackage{amsmath, amssymb, amsthm}
\usepackage{booktabs, array, multirow, tabularx}
\usepackage{xcolor}
\usepackage{tcolorbox}
\usepackage{hyperref}
\usepackage{enumitem}

\tcbuselibrary{breakable}

\title{Appendix XIX: LIT-LEARNING\\
Experience Goods, Belief Updating, and Learning in Behavioral Economics}
\author{BEATRIX Research Group \\ FehrAdvice \& Partners AG}
\date{January 2026}

\begin{document}

\maketitle

%===============================================================================
% METADATA & OVERVIEW
%===============================================================================

\begin{tcolorbox}[colback=gray!10!white, colframe=gray!75!black, title=Appendix Metadata]
\begin{tabular}{@{}ll@{}}
\textbf{Code:} & XIX (LIT-LEARNING) \\
\textbf{Category:} & LIT: Literature Integration \\
\textbf{Version:} & 1.0 (January 2026) \\
\textbf{Purpose:} & Consolidate learning and belief updating literature for behavioral economics \\
\textbf{Related Chapters:} & Ch. 5 (Awareness), Ch. 9 (Context), Ch. 14 (Interventions) \\
\textbf{Related COREs:} & AU (CORE-AWARE), AW (CORE-STAGE), BBB (CORE-WHERE) \\
\textbf{Related Appendices:} & HF (PREDICT-HABIT), PBB (FORMAL-PROBABILITY) \\
\end{tabular}
\end{tcolorbox}

%===============================================================================
% CROSS-REFERENCE MAP
%===============================================================================

\begin{tcolorbox}[colback=cyan!5!white, colframe=cyan!75!black, title=Cross-Reference Map]
\textbf{Incoming References (Who cites XIX):}
\begin{itemize}[nosep]
    \item \textbf{Appendix HF (PREDICT-HABIT):} Theoretical foundation for experience-good learning
    \item \textbf{Appendix UNMAPPED_PBB (FORMAL-PROBABILITY):} Empirical grounding for Axiom P7 (Learning Feedback)
    \item \textbf{Appendix UNMAPPED_AU (CORE-AWARE):} Literature on awareness and belief updating
    \item \textbf{Chapter 5:} Awareness function derivation
\end{itemize}

\textbf{Outgoing References (What XIX depends on):}
\begin{itemize}[nosep]
    \item \textbf{Master Bibliography:} bcm\_master.bib (Section: APPENDIX HF)
    \item \textbf{Appendix G (Glossary):} Definitions of learning-related terms
    \item \textbf{Appendix UNMAPPED_BBB (CORE-WHERE):} Parameter estimates from learning studies
\end{itemize}
\end{tcolorbox}

%===============================================================================
% CHAPTER LINKAGE
%===============================================================================

\begin{tcolorbox}[colback=blue!5!white, colframe=blue!75!black, title=Chapter Linkage]
\textbf{Primary Chapter:} Chapter 5 (The Awareness Function)
\begin{itemize}[nosep]
    \item Learning literature grounds the awareness updating mechanism $A(\cdot)$
    \item Experience goods explain why $\Delta A^{(1)} \gg \Delta A^{(k)}$ for $k \geq 2$
\end{itemize}

\textbf{Secondary Chapters:}
\begin{itemize}[nosep]
    \item Chapter 9 (Context): Learning as context-dependent process
    \item Chapter 14 (Interventions): Front-loading principle for intervention design
    \item Chapter 16 (Probability): Learning feedback in behavioral change probability
\end{itemize}
\end{tcolorbox}

%===============================================================================
% QUICK REFERENCE
%===============================================================================

\begin{tcolorbox}[colback=yellow!5!white, colframe=yellow!75!black, title=Quick Reference]
\textbf{Core Question:} What does behavioral economics know about learning, belief updating, and experience goods---and how does this inform intervention design?

\textbf{Key Literature Streams:}
\begin{enumerate}[nosep]
    \item \textbf{Experience Goods:} Nelson (1970), Ackerberg (2003)---quality learned through consumption
    \item \textbf{Bayesian Learning:} Rational belief updating from signals
    \item \textbf{Habit vs. Addiction:} Becker \& Murphy (1988) vs. Jones et al. (2025)
    \item \textbf{Development RCTs:} Dupas (2014), Banerjee et al. (2021)---learning in low-income settings
    \item \textbf{Incentive Persistence:} Charness \& Gneezy (2009), Royer et al. (2015)
\end{enumerate}

\textbf{Central Insight:} For experience goods, learning is a \textit{discrete transition}, not a \textit{gradual accumulation}. First exposure dominates subsequent exposures.

\textbf{EBF Integration:} Grounds Axiom P7 (Learning Feedback), informs CORE-AWARE ($A(\cdot)$), and provides parameters for CORE-WHERE ($\Theta$).
\end{tcolorbox}

%===============================================================================
% ABSTRACT
%===============================================================================

\begin{abstract}
This appendix consolidates the behavioral economics literature on learning, belief updating, and experience goods. We distinguish four major literature streams: (1) experience goods theory (Nelson 1970, Ackerberg 2003), (2) addiction vs. learning models (Becker \& Murphy 1988 vs. Jones et al. 2025), (3) development RCTs on learning (Dupas 2014, Banerjee et al. 2021), and (4) incentive persistence studies (Charness \& Gneezy 2009, Allcott \& Rogers 2014). The central finding is that for experience goods, learning operates through a \textit{discrete transition} at first exposure rather than gradual accumulation. This has profound implications for intervention design: front-load investment in first experience quality, as subsequent interventions have diminishing returns. The appendix provides parameter estimates for CORE-WHERE and grounds key EBF axioms including UNMAPPED_AU-3 (First Experience Dominance) and P7 (Learning Feedback).
\end{abstract}

%===============================================================================
% FUNDAMENTAL QUESTION
%===============================================================================

% -----------------------------------------------------------------------------
% APPENDIX SCOPE BOX
% -----------------------------------------------------------------------------

\begin{tcolorbox}[
    colback=orange!5!white,
    colframe=orange!75!black,
    title={\textbf{Appendix Scope: Ziel / In-Scope / Out-of-Scope / Constraints}},
    fonttitle=\bfseries
]

\textbf{Ziel (Objective):}
\begin{itemize}[nosep]
    \item Classic Economics:
\end{itemize}

\textbf{In-Scope:}
\begin{itemize}[nosep]
    \item Fundamental Question
    \item Literature Stream 1: Experience Goods
    \item Literature Stream 2: Addiction vs. Learning Models
    \item Literature Stream 3: Learning in Development Economics
    \item Literature Stream 4: Incentive Persistence
\end{itemize}

\textbf{Out-of-Scope (delegiert an andere Appendices):}
\begin{itemize}[nosep]
    \item PREDICT-HABIT $\rightarrow$ Appendix HF
    \item FORMAL-PROBABILITY $\rightarrow$ Appendix UNMAPPED_PBB
    \item CORE-AWARE $\rightarrow$ Appendix UNMAPPED_AU
    \item Glossary $\rightarrow$ Appendix G
    \item CORE-WHERE $\rightarrow$ Appendix UNMAPPED_BBB
\end{itemize}

\textbf{Constraints (Anwendungsgrenzen):}
\begin{itemize}[nosep]
    \item [TODO: Define when this appendix does NOT apply]
    \item [TODO: Define required prerequisites]
    \item [TODO: Define known limitations]
\end{itemize}

\textbf{Lieferobjekte (Deliverables):}
\begin{itemize}[nosep]
    \item 5 Table(s)
    \item 10 Equation(s)
    \item 3 Axiom(s)
\end{itemize}

\end{tcolorbox}


\section{Fundamental Question}
\label{sec:lrn-question}

\begin{tcolorbox}[colback=orange!5!white, colframe=orange!75!black,
    title=The Learning Question in Behavioral Economics]

\textbf{Classic Economics:} Agents have stable, known preferences. Learning is about \textit{information} (prices, probabilities), not \textit{utility}.

\textbf{Behavioral Economics Challenge:} Agents often don't know their own utility until they experience an option. Learning is about \textit{discovering preferences}, not just \textit{updating beliefs about states}.

\textbf{Central Question:}
\begin{quote}
\textit{How do agents learn about the value of goods and behaviors they have never experienced---and what are the dynamics of this learning process?}
\end{quote}

\textbf{Why This Matters for Interventions:}
\begin{itemize}[nosep]
    \item If learning is \textbf{gradual} $\Rightarrow$ sustained intervention needed
    \item If learning is \textbf{one-shot} $\Rightarrow$ front-load investment in first experience
    \item If learning is \textbf{heterogeneous} $\Rightarrow$ segment-specific strategies required
\end{itemize}

\end{tcolorbox}


%%%%%%%%%%%%%%%%%%%%%%%%%%%%%%%%%%%%%%%%%%%%%%%%%%%%%%%%%%%%%%%%%%%%%%%%%%%%%%%
\section{Literature Stream 1: Experience Goods}
\label{sec:lrn-experience-goods}
%%%%%%%%%%%%%%%%%%%%%%%%%%%%%%%%%%%%%%%%%%%%%%%%%%%%%%%%%%%%%%%%%%%%%%%%%%%%%%%

\subsection{Foundational Theory: Nelson (1970)}

Philip Nelson introduced the distinction between \textbf{search goods} and \textbf{experience goods}:

\begin{table}[htbp]
\centering
\caption{Search vs. Experience Goods}
\label{tab:lrn-search-experience}
\renewcommand{\arraystretch}{1.3}
\begin{tabular}{@{}lll@{}}
\toprule
\textbf{Property} & \textbf{Search Goods} & \textbf{Experience Goods} \\
\midrule
Quality observable & Before purchase & Only after consumption \\
Information source & Inspection, advertising & Trial, reputation, reviews \\
Examples & Clothing fit, car specs & Restaurant quality, health screening \\
Learning timing & Pre-decision & Post-decision \\
\bottomrule
\end{tabular}
\end{table}

\textbf{Key Insight:} For experience goods, agents \textit{cannot} fully evaluate utility before trying. This creates systematic \textbf{undervaluation} of unfamiliar options.

\subsection{Empirical Identification: Ackerberg (2003)}

Daniel Ackerberg provided the first clean empirical identification of experience-good effects:

\begin{tcolorbox}[colback=gray!5!white, colframe=gray!75!black,
    title=Ackerberg (2003): Advertising and Experience Goods]

\textbf{Setting:} New yogurt brand introduction

\textbf{Identification Strategy:} Compare advertising effects for:
\begin{itemize}[nosep]
    \item Consumers who have \textit{not} tried the product (experience effect possible)
    \item Consumers who \textit{have} tried the product (experience effect resolved)
\end{itemize}

\textbf{Key Finding:} Advertising affects \textit{inexperienced} consumers but \textit{not} experienced ones.

\textbf{Interpretation:} Advertising provides information about experience goods that is \textit{dominated} by direct experience once available.

\textbf{Implication for EBF:} First trial resolves uncertainty; subsequent information (advertising, nudges) has diminishing returns.

\end{tcolorbox}

\subsection{EBF Integration}

\begin{tcolorbox}[colback=green!5!white, colframe=green!75!black,
    title=Experience Goods $\rightarrow$ CORE-AWARE]

The experience goods literature grounds the EBF awareness function $A(\cdot)$:

\textbf{Pre-experience:}
\begin{equation}
A_0 = A(\text{advertising}, \text{social proof}, \text{defaults}) < A_{\text{true}}
\end{equation}

\textbf{Post-experience:}
\begin{equation}
A_1 = A(\text{direct experience}) \approx A_{\text{true}}
\end{equation}

\textbf{Key Property:} $\Delta A^{(1)} = A_1 - A_0 \gg \Delta A^{(k)} \approx 0$ for $k \geq 2$

This is \textbf{Axiom UNMAPPED_AU-3} in CORE-AWARE: ``First experience dominates subsequent information.''

\end{tcolorbox}


%%%%%%%%%%%%%%%%%%%%%%%%%%%%%%%%%%%%%%%%%%%%%%%%%%%%%%%%%%%%%%%%%%%%%%%%%%%%%%%
\section{Literature Stream 2: Addiction vs. Learning Models}
\label{sec:lrn-addiction-vs-learning}
%%%%%%%%%%%%%%%%%%%%%%%%%%%%%%%%%%%%%%%%%%%%%%%%%%%%%%%%%%%%%%%%%%%%%%%%%%%%%%%

\subsection{The Addiction Paradigm: Becker \& Murphy (1988)}

The canonical model of habit formation in economics is \textbf{rational addiction}:

\begin{tcolorbox}[colback=red!5!white, colframe=red!75!black,
    title=Becker \& Murphy (1988): Rational Addiction Model]

\textbf{Core Mechanism:} Past consumption increases current marginal utility of consumption (adjacent complementarity).

\textbf{State Variable:} Consumption stock $S_t$ accumulates:
\begin{equation}
S_{t+1} = (1-\delta) S_t + c_t
\end{equation}
where $\delta$ is the depreciation rate.

\textbf{Utility:}
\begin{equation}
U(c_t, S_t) = u(c_t) + \alpha \cdot c_t \cdot S_t
\end{equation}
where $\alpha > 0$ captures reinforcement.

\textbf{Key Prediction:} \textit{Recent} consumption matters more than \textit{distant} consumption:
\begin{equation}
\frac{\partial c_3}{\partial c_2} > \frac{\partial c_3}{\partial c_1}
\end{equation}

\textbf{Domains Where Model Fits:}
\begin{itemize}[nosep]
    \item Substance addiction (cigarettes, alcohol, drugs)
    \item High-frequency consumption goods
    \item Goods with physiological reinforcement
\end{itemize}

\end{tcolorbox}

\subsection{The Learning Alternative: Jones, Molitor \& Reif (2025)}

For many behaviors, the addiction model is \textbf{falsified}:

\begin{tcolorbox}[colback=blue!5!white, colframe=blue!75!black,
    title=Jones, Molitor \& Reif (2025): Experience-Good Learning]

\textbf{Setting:} Annual biometric health screenings (N = 4,799 employees)

\textbf{Design Innovation:} Rerandomization of incentives each year allows separate identification of $\theta_1$ (first exposure effect) and $\theta_2$ (second exposure effect).

\textbf{Key Finding:}
\begin{equation}
\theta_1 = 0.32\text{--}0.36 \quad \text{vs.} \quad \theta_2 \approx 0.03 \text{ (n.s.)}
\end{equation}

\textbf{Interpretation:} First exposure has large, permanent effect; second exposure adds almost nothing.

\textbf{Addiction Model Prediction:} $\theta_2 > \theta_1$ (recent dominates)

\textbf{Verdict:} Addiction model \textbf{rejected}; experience-good learning \textbf{supported}.

\end{tcolorbox}

\subsection{Mechanism Discrimination Framework}

\begin{table}[htbp]
\centering
\caption{When Does Each Model Apply?}
\label{tab:lrn-mechanism-domains}
\renewcommand{\arraystretch}{1.4}
\begin{tabular}{@{}p{3cm}p{4cm}p{4.5cm}@{}}
\toprule
\textbf{Model} & \textbf{Testable Prediction} & \textbf{Domain Examples} \\
\midrule
\textbf{Addiction} & Recent $>$ distant: $\theta_2 > \theta_1$ & Cigarettes, alcohol, gambling, social media \\
\textbf{Learning} & First $\gg$ subsequent: $\theta_1 \gg \theta_2 \approx 0$ & Health screenings, retirement savings, sustainable behaviors \\
\textbf{Switching Costs} & Any prior sufficient: $\exists$ matters, not timing & Software platforms, bank accounts \\
\bottomrule
\end{tabular}
\end{table}

\textbf{EBF Implication:} The mechanism determines the optimal intervention strategy. See Appendix HF (PREDICT-HABIT) for formal taxonomy.


%%%%%%%%%%%%%%%%%%%%%%%%%%%%%%%%%%%%%%%%%%%%%%%%%%%%%%%%%%%%%%%%%%%%%%%%%%%%%%%
\section{Literature Stream 3: Learning in Development Economics}
\label{sec:lrn-development}
%%%%%%%%%%%%%%%%%%%%%%%%%%%%%%%%%%%%%%%%%%%%%%%%%%%%%%%%%%%%%%%%%%%%%%%%%%%%%%%

Development economics has produced some of the cleanest evidence on learning from experience, often using RCT designs.

\subsection{Dupas (2014): Learning About Health Products}

\begin{tcolorbox}[colback=gray!5!white, colframe=gray!75!black,
    title=Dupas (2014): Short-Run Subsidies and Long-Run Adoption]

\textbf{Context:} Insecticide-treated bednets in Kenya

\textbf{Question:} Do short-run subsidies crowd out long-run willingness to pay?

\textbf{Design:} Randomize initial subsidy levels, then offer repurchase at positive price.

\textbf{Key Finding:} Initial subsidies \textit{increase} subsequent willingness to pay.

\textbf{Mechanism:} Experience with product reveals quality previously unknown.

\textbf{Implication:} For experience goods, subsidizing first trial can be \textit{welfare-improving} even if subsidies are later removed.

\end{tcolorbox}

\subsection{Banerjee et al. (2021): Choosing to Exercise}

\begin{tcolorbox}[colback=gray!5!white, colframe=gray!75!black,
    title=Banerjee et al. (2021): Exercise Habits and Incentives]

\textbf{Context:} Gym attendance among Indian office workers

\textbf{Design:} Randomize incentive timing and structure

\textbf{Key Finding:} Incentives work primarily through \textit{learning about preferences}, not through forming ``habits'' in the addiction sense.

\textbf{Evidence:}
\begin{itemize}[nosep]
    \item Post-incentive attendance predicts reported enjoyment
    \item Effect strongest for those with initial uncertainty about preferences
    \item No dose-response relationship (more visits $\not\Rightarrow$ stronger habit)
\end{itemize}

\textbf{Implication:} Gym ``habits'' are better understood as resolved preference uncertainty than as accumulated consumption capital.

\end{tcolorbox}

\subsection{Hussam et al. (2022): Habit Formation in Handwashing}

\begin{tcolorbox}[colback=gray!5!white, colframe=gray!75!black,
    title=Hussam et al. (2022): Targeting High Frequency Behaviors]

\textbf{Context:} Handwashing behavior in Bangladesh

\textbf{Design:} Multi-arm RCT with different intervention intensities

\textbf{Key Finding:} For \textit{high-frequency} behaviors (handwashing multiple times per day), the addiction/reinforcement model \textit{does} fit better.

\textbf{Contrast with Jones et al.:}
\begin{itemize}[nosep]
    \item Health screenings: Annual, low-frequency $\Rightarrow$ Learning model
    \item Handwashing: Daily, high-frequency $\Rightarrow$ Reinforcement model
\end{itemize}

\textbf{Implication:} Behavior frequency moderates which mechanism dominates.

\end{tcolorbox}


%%%%%%%%%%%%%%%%%%%%%%%%%%%%%%%%%%%%%%%%%%%%%%%%%%%%%%%%%%%%%%%%%%%%%%%%%%%%%%%
\section{Literature Stream 4: Incentive Persistence}
\label{sec:lrn-incentive-persistence}
%%%%%%%%%%%%%%%%%%%%%%%%%%%%%%%%%%%%%%%%%%%%%%%%%%%%%%%%%%%%%%%%%%%%%%%%%%%%%%%

A key question for intervention design: Do incentive effects persist after incentives are removed?

\subsection{Charness \& Gneezy (2009): Gym Attendance}

\begin{tcolorbox}[colback=gray!5!white, colframe=gray!75!black,
    title=Charness \& Gneezy (2009): Exercise and Incentives]

\textbf{Setting:} University gym attendance

\textbf{Design:} Pay students to attend gym for 4 weeks, then observe post-incentive behavior

\textbf{Key Finding:} Incentivized attendance \textit{persists} after incentives removed.

\textbf{Interpretation:} Incentives helped students \textit{discover} they enjoyed exercise more than expected.

\textbf{Limitation:} Cannot distinguish learning from habit formation with this design alone.

\end{tcolorbox}

\subsection{Royer, Stehr \& Sydnor (2015): Commitment Devices}

\begin{tcolorbox}[colback=gray!5!white, colframe=gray!75!black,
    title=Royer, Stehr \& Sydnor (2015): Incentives and Commitment]

\textbf{Setting:} Workplace wellness program

\textbf{Design:} Compare:
\begin{itemize}[nosep]
    \item Simple incentives (payment for gym visits)
    \item Commitment contracts (deposit own money, forfeit if miss target)
\end{itemize}

\textbf{Key Finding:} Commitment contracts show \textit{stronger} persistence than simple incentives.

\textbf{Interpretation:} Commitment devices may trigger deeper engagement, leading to more learning.

\textbf{EBF Connection:} Maps to CORE-READY (AV)---commitment changes willingness threshold $\theta$.

\end{tcolorbox}

\subsection{Allcott \& Rogers (2014): Energy Conservation}

\begin{tcolorbox}[colback=gray!5!white, colframe=gray!75!black,
    title=Allcott \& Rogers (2014): Persistence of Social Norm Messages]

\textbf{Setting:} Opower home energy reports (social comparison nudges)

\textbf{Design:} Observe behavior during and after treatment cessation

\textbf{Key Finding:} Effects decay after treatment stops, but \textit{not} to zero.

\textbf{Decay Pattern:}
\begin{itemize}[nosep]
    \item During treatment: 2\% energy reduction
    \item 2 years post-treatment: Still 1\% reduction (50\% persistence)
\end{itemize}

\textbf{Interpretation:} Partial learning effect---some households permanently update beliefs about energy use.

\textbf{EBF Implication:} Learning-based persistence may be partial, not all-or-nothing.

\end{tcolorbox}


%%%%%%%%%%%%%%%%%%%%%%%%%%%%%%%%%%%%%%%%%%%%%%%%%%%%%%%%%%%%%%%%%%%%%%%%%%%%%%%
\section{Parameter Estimates for CORE-WHERE}
\label{sec:lrn-parameters}
%%%%%%%%%%%%%%%%%%%%%%%%%%%%%%%%%%%%%%%%%%%%%%%%%%%%%%%%%%%%%%%%%%%%%%%%%%%%%%%

The learning literature provides empirical estimates for EBF parameter repository (Appendix UNMAPPED_BBB):

\begin{table}[htbp]
\centering
\caption{Learning Parameters from Empirical Literature}
\label{tab:lrn-parameters}
\renewcommand{\arraystretch}{1.4}
\begin{tabular}{@{}llcl@{}}
\toprule
\textbf{Parameter} & \textbf{Description} & \textbf{Estimate} & \textbf{Source} \\
\midrule
$\theta_1$ & First exposure effect (health) & 0.32--0.36 & Jones et al. (2025) \\
$\theta_2$ & Second exposure effect (health) & $\approx 0.03$ (n.s.) & Jones et al. (2025) \\
$\theta_1/\theta_2$ & Learning asymmetry ratio & $> 10$ & Jones et al. (2025) \\
Persistence & Post-incentive retention rate & 0.50 & Allcott \& Rogers (2014) \\
Ad effect (inexperienced) & Advertising on non-triers & Significant & Ackerberg (2003) \\
Ad effect (experienced) & Advertising on triers & $\approx 0$ & Ackerberg (2003) \\
Subsidy spillover & Effect on later WTP & Positive & Dupas (2014) \\
\bottomrule
\end{tabular}
\end{table}

\textbf{Key Calibration Insight:} For low-frequency experience goods, assume $\theta_1 \approx 0.35$ and $\theta_2 \approx 0$ as default priors.


%%%%%%%%%%%%%%%%%%%%%%%%%%%%%%%%%%%%%%%%%%%%%%%%%%%%%%%%%%%%%%%%%%%%%%%%%%%%%%%
\section{Theoretical Integration with EBF}
\label{sec:lrn-integration}
%%%%%%%%%%%%%%%%%%%%%%%%%%%%%%%%%%%%%%%%%%%%%%%%%%%%%%%%%%%%%%%%%%%%%%%%%%%%%%%

\subsection{Connection to CORE-AWARE (AU)}

\begin{tcolorbox}[colback=blue!5!white, colframe=blue!75!black,
    title=Integration: LIT-LEARNING $\leftrightarrow$ CORE-AWARE]

The learning literature provides the \textbf{micro-foundation} for the awareness function $A(\cdot)$:

\textbf{Axiom UNMAPPED_AU-3 (First Experience Dominance):}
\begin{equation}
\Delta A^{(1)} \gg \Delta A^{(k)} \quad \text{for } k \geq 2
\end{equation}

\textbf{Empirical Support:}
\begin{itemize}[nosep]
    \item Ackerberg (2003): Advertising effect vanishes post-experience
    \item Jones et al. (2025): $\theta_1 \gg \theta_2$
    \item Banerjee et al. (2021): Preference discovery, not accumulation
\end{itemize}

\end{tcolorbox}

\subsection{Connection to FORMAL-PROBABILITY (PBB)}

\begin{tcolorbox}[colback=blue!5!white, colframe=blue!75!black,
    title=Integration: LIT-LEARNING $\leftrightarrow$ FORMAL-PROBABILITY]

The learning literature grounds \textbf{Axiom P7 (Learning Feedback)} and \textbf{Theorem T3 (Learning Reduces Barriers)}:

\textbf{Axiom P7:} Probabilities adapt through experience:
\begin{equation}
P(i \to j, t+\Delta t) = f(P(i \to j, t), \text{Experience}(t))
\end{equation}

\textbf{Theorem T3:} Repeated exposure reduces effective threshold:
\begin{equation}
\theta_{\text{eff}}(k) = \theta_0 \cdot g(k), \quad g'(k) < 0
\end{equation}

\textbf{Refinement from Learning Literature:}
For experience goods, $g(k)$ is approximately a \textit{step function}:
\begin{equation}
g(k) \approx \begin{cases}
1 & \text{if } k = 0 \\
g_1 \ll 1 & \text{if } k \geq 1
\end{cases}
\end{equation}

\end{tcolorbox}

\subsection{Connection to PREDICT-HABIT (HF)}

\begin{tcolorbox}[colback=blue!5!white, colframe=blue!75!black,
    title=Integration: LIT-LEARNING $\leftrightarrow$ PREDICT-HABIT]

Appendix HF (PREDICT-HABIT) uses this literature to:

\begin{enumerate}[nosep]
    \item \textbf{Ground the mechanism taxonomy:} Addiction vs. Learning vs. Switching Costs
    \item \textbf{Provide empirical calibration:} $\theta_1, \theta_2$ from Jones et al.
    \item \textbf{Derive the absorbing-state model:} $H_{i,t+1} = \max(H_{it}, S_{it})$
    \item \textbf{Support the front-loading principle:} 70\% ROI gain from first-exposure investment
\end{enumerate}

\textbf{Direction:} LIT-LEARNING (this appendix) provides literature; PREDICT-HABIT provides predictions.

\end{tcolorbox}


%%%%%%%%%%%%%%%%%%%%%%%%%%%%%%%%%%%%%%%%%%%%%%%%%%%%%%%%%%%%%%%%%%%%%%%%%%%%%%%
\section{Policy Implications}
\label{sec:lrn-policy}
%%%%%%%%%%%%%%%%%%%%%%%%%%%%%%%%%%%%%%%%%%%%%%%%%%%%%%%%%%%%%%%%%%%%%%%%%%%%%%%

\begin{tcolorbox}[colback=green!10!white, colframe=green!75!black,
    title=Policy Design Principles from Learning Literature]

\textbf{Principle 1: Front-Load Investment}
\begin{quote}
For experience goods, invest heavily in \textit{first} participation quality. Subsequent nudges have diminishing returns.
\end{quote}

\textbf{Principle 2: Subsidize Trial, Not Repetition}
\begin{quote}
Temporary subsidies for first adoption can be welfare-improving if they trigger learning. Ongoing subsidies may be wasteful.
\end{quote}

\textbf{Principle 3: Match Mechanism to Behavior Frequency}
\begin{quote}
\begin{itemize}[nosep]
    \item Low-frequency behaviors (annual events): Use learning model
    \item High-frequency behaviors (daily habits): Consider reinforcement model
\end{itemize}
\end{quote}

\textbf{Principle 4: Persistence Varies by Domain}
\begin{quote}
Expect 50--100\% persistence for true experience goods (health, savings). Expect lower persistence for behaviors without clear quality revelation.
\end{quote}

\end{tcolorbox}


%%%%%%%%%%%%%%%%%%%%%%%%%%%%%%%%%%%%%%%%%%%%%%%%%%%%%%%%%%%%%%%%%%%%%%%%%%%%%%%
\section{Summary}
\label{sec:lrn-summary}
%%%%%%%%%%%%%%%%%%%%%%%%%%%%%%%%%%%%%%%%%%%%%%%%%%%%%%%%%%%%%%%%%%%%%%%%%%%%%%%

\begin{tcolorbox}[colback=green!10!white, colframe=green!75!black,
    title=\textbf{Appendix XIX: Summary}]

\textbf{Core Question:} How do agents learn about the value of unfamiliar goods and behaviors?

\textbf{Key Literature Streams:}
\begin{enumerate}[nosep]
    \item \textbf{Experience Goods:} Nelson (1970), Ackerberg (2003)
    \item \textbf{Addiction vs. Learning:} Becker \& Murphy (1988) vs. Jones et al. (2025)
    \item \textbf{Development RCTs:} Dupas (2014), Banerjee et al. (2021), Hussam et al. (2022)
    \item \textbf{Incentive Persistence:} Charness \& Gneezy (2009), Royer et al. (2015), Allcott \& Rogers (2014)
\end{enumerate}

\textbf{Central Insight:}
\begin{quote}
\textit{For experience goods, learning is a discrete transition at first exposure, not a gradual accumulation. This has profound implications for intervention design: front-load investment in first experience quality.}
\end{quote}

\textbf{EBF Integration:}
\begin{itemize}[nosep]
    \item Grounds CORE-AWARE (AU): Axiom UNMAPPED_AU-3 (First Experience Dominance)
    \item Supports FORMAL-PROBABILITY (PBB): Axiom P7, Theorem T3
    \item Provides parameters for CORE-WHERE (BBB): $\theta_1 \approx 0.35$, persistence $\approx 50\%$
    \item Informs PREDICT-HABIT (HF): Mechanism discrimination framework
\end{itemize}

\end{tcolorbox}


%%%%%%%%%%%%%%%%%%%%%%%%%%%%%%%%%%%%%%%%%%%%%%%%%%%%%%%%%%%%%%%%%%%%%%%%%%%%%%%
\section{Annotated Bibliography}
\label{sec:lrn-bibliography}
%%%%%%%%%%%%%%%%%%%%%%%%%%%%%%%%%%%%%%%%%%%%%%%%%%%%%%%%%%%%%%%%%%%%%%%%%%%%%%%

\subsection{Foundational Papers}

\begin{enumerate}

\item \textbf{Nelson, P. (1970).} ``Information and Consumer Behavior.'' \textit{Journal of Political Economy}.
\begin{quote}
\textit{Introduces search vs. experience goods distinction. Foundational for understanding why agents undervalue unfamiliar options.}
\end{quote}

\item \textbf{Becker, G. S. \& Murphy, K. M. (1988).} ``A Theory of Rational Addiction.'' \textit{Journal of Political Economy}.
\begin{quote}
\textit{The canonical addiction model. Predicts recent consumption dominates. Falsified for experience goods by Jones et al. (2025).}
\end{quote}

\item \textbf{Ackerberg, D. (2003).} ``Advertising, Learning, and Consumer Choice in Experience Good Markets.'' \textit{International Economic Review}.
\begin{quote}
\textit{First clean identification of experience-good effects. Advertising affects inexperienced but not experienced consumers.}
\end{quote}

\end{enumerate}

\subsection{Key Empirical Papers}

\begin{enumerate}[resume]

\item \textbf{Dupas, P. (2014).} ``Short-Run Subsidies and Long-Run Adoption of New Health Products.'' \textit{Econometrica}.
\begin{quote}
\textit{Shows subsidies for experience goods can increase later WTP. Learning dominates crowding-out.}
\end{quote}

\item \textbf{Jones, D., Molitor, D. \& Reif, J. (2025).} ``What Do Workplace Wellness Programs Do?'' \textit{Quarterly Journal of Economics}.
\begin{quote}
\textit{The key paper for PREDICT-HABIT. Cleanly identifies $\theta_1 \gg \theta_2$ using rerandomization design. Falsifies addiction model for health screenings.}
\end{quote}

\item \textbf{Banerjee, A. et al. (2021).} ``Choosing to Exercise.'' \textit{AER: Insights}.
\begin{quote}
\textit{Demonstrates gym ``habits'' are preference discovery, not accumulated consumption capital.}
\end{quote}

\item \textbf{Hussam, R. et al. (2022).} ``Habit Formation and Rational Addiction.'' \textit{Journal of Political Economy}.
\begin{quote}
\textit{Shows reinforcement model fits better for high-frequency behaviors (handwashing). Provides boundary condition for learning model.}
\end{quote}

\item \textbf{Bregolin, J., Hopfensitz, A. \& Panova, E. (2026).} ``Using Advice Without Considering Its Quality: A Laboratory Experiment of Demand for Advice.'' \textit{Working Paper}.
\begin{quote}
\textit{Tests reputational cheap-talk theory. Key finding: $\sim$50\% of participants do NOT update beliefs about advisor quality. Demand for advice is independent of confirmatory vs. contrarian alignment. Hot-hand fallacy drives purchasing behavior. Challenges Bayesian learning assumptions.}
\end{quote}

\end{enumerate}

\subsection{Methodological Papers}

\begin{enumerate}[resume]

\item \textbf{Charness, G. \& Gneezy, U. (2009).} ``Incentives to Exercise.'' \textit{Econometrica}.
\begin{quote}
\textit{Early evidence of incentive persistence. Cannot distinguish mechanisms but motivates subsequent work.}
\end{quote}

\item \textbf{Royer, H., Stehr, M. \& Sydnor, J. (2015).} ``Incentives, Commitments, and Habit Formation.'' \textit{AER}.
\begin{quote}
\textit{Compares incentive structures. Commitment contracts show stronger persistence than simple payments.}
\end{quote}

\item \textbf{Allcott, H. \& Rogers, T. (2014).} ``The Short-Run and Long-Run Effects of Behavioral Interventions.'' \textit{AER}.
\begin{quote}
\textit{Documents partial persistence (50\%) of social norm interventions. Provides realistic calibration for policy.}
\end{quote}

\end{enumerate}


%%%%%%%%%%%%%%%%%%%%%%%%%%%%%%%%%%%%%%%%%%%%%%%%%%%%%%%%%%%%%%%%%%%%%%%%%%%%%%%
\section{Key Results}
\label{sec:lrn-results}
%%%%%%%%%%%%%%%%%%%%%%%%%%%%%%%%%%%%%%%%%%%%%%%%%%%%%%%%%%%%%%%%%%%%%%%%%%%%%%%

\begin{tcolorbox}[colback=purple!5!white, colframe=purple!75!black,
    title=Summary of Empirical Results from Learning Literature]

\textbf{Result 1: First Exposure Dominance (Jones et al. 2025)}
\begin{itemize}[nosep]
    \item $\theta_1 = 0.32$--$0.36$ (first exposure effect)
    \item $\theta_2 \approx 0.03$ (n.s.) (second exposure effect)
    \item Ratio $\theta_1/\theta_2 > 10$
\end{itemize}

\textbf{Result 2: Advertising-Experience Asymmetry (Ackerberg 2003)}
\begin{itemize}[nosep]
    \item Advertising affects inexperienced consumers: Significant
    \item Advertising affects experienced consumers: $\approx 0$
\end{itemize}

\textbf{Result 3: Incentive Persistence (Allcott \& Rogers 2014)}
\begin{itemize}[nosep]
    \item During treatment: 2\% energy reduction
    \item 2 years post-treatment: 1\% reduction (50\% persistence)
\end{itemize}

\textbf{Result 4: Frequency Moderates Mechanism (Hussam et al. 2022)}
\begin{itemize}[nosep]
    \item Low-frequency behaviors: Learning model fits better
    \item High-frequency behaviors: Reinforcement model fits better
\end{itemize}

\end{tcolorbox}


%%%%%%%%%%%%%%%%%%%%%%%%%%%%%%%%%%%%%%%%%%%%%%%%%%%%%%%%%%%%%%%%%%%%%%%%%%%%%%%
\section{Glossary of Terms}
\label{sec:lrn-glossary}
%%%%%%%%%%%%%%%%%%%%%%%%%%%%%%%%%%%%%%%%%%%%%%%%%%%%%%%%%%%%%%%%%%%%%%%%%%%%%%%

\textit{For complete definitions, see Appendix G (REF-GLOSSARY).}

\begin{table}[htbp]
\centering
\caption{Key Terms in Learning Literature}
\label{tab:lrn-glossary}
\renewcommand{\arraystretch}{1.3}
\begin{tabular}{@{}lp{8cm}c@{}}
\toprule
\textbf{Term} & \textbf{Definition} & \textbf{In G?} \\
\midrule
Experience good & Good whose quality can only be evaluated after consumption & NEW \\
Search good & Good whose quality can be evaluated before purchase & NEW \\
Belief updating & Bayesian revision of probability distributions given new information & \checkmark \\
Absorbing state & State from which no further transitions occur & NEW \\
First exposure effect ($\theta_1$) & Causal impact of first participation on subsequent behavior & NEW \\
Incentive persistence & Continued behavior change after incentive removal & NEW \\
Preference discovery & Learning about own utility through experience & NEW \\
Reinforcement & Mechanism where past consumption increases current marginal utility & \checkmark \\
\bottomrule
\end{tabular}
\end{table}


%%%%%%%%%%%%%%%%%%%%%%%%%%%%%%%%%%%%%%%%%%%%%%%%%%%%%%%%%%%%%%%%%%%%%%%%%%%%%%%
\section{Critical Foundations and Limitations}
\label{sec:lrn-foundations}
%%%%%%%%%%%%%%%%%%%%%%%%%%%%%%%%%%%%%%%%%%%%%%%%%%%%%%%%%%%%%%%%%%%%%%%%%%%%%%%

\subsection{Foundation 1: Rational Learning Assumption}

\textbf{Assumption:} Agents update beliefs according to Bayes' rule given experience signals.

\textbf{Support:} Jones et al. (2025) results are consistent with rational learning; no evidence of systematic updating biases.

\textbf{Limitation:} Does not rule out bounded rationality in other dimensions (attention, memory).

\subsection{Foundation 2: Signal Quality Homogeneity}

\textbf{Assumption:} First experience provides a representative signal of true quality.

\textbf{Support:} Absorbing-state model predicts $\theta_2 \approx 0$, which is observed.

\textbf{Limitation:} Does not identify effect of \textit{quality} of first experience (good vs. bad). See PREDICT-HABIT Identification Boundary.

\subsection{Foundation 3: Domain Generalizability}

\textbf{Assumption:} Learning model applies broadly to experience goods.

\textbf{Support:} Consistent evidence across health (Jones et al.), consumer goods (Ackerberg), development (Dupas).

\textbf{Limitation:} Frequency moderates mechanism (Hussam et al. 2022)---high-frequency behaviors may follow reinforcement model.


%%%%%%%%%%%%%%%%%%%%%%%%%%%%%%%%%%%%%%%%%%%%%%%%%%%%%%%%%%%%%%%%%%%%%%%%%%%%%%%
\section{Open Issues and Research Agenda}
\label{sec:lrn-open}
%%%%%%%%%%%%%%%%%%%%%%%%%%%%%%%%%%%%%%%%%%%%%%%%%%%%%%%%%%%%%%%%%%%%%%%%%%%%%%%

\begin{table}[htbp]
\centering
\caption{Open Questions in Learning Literature}
\label{tab:lrn-open}
\renewcommand{\arraystretch}{1.3}
\begin{tabular}{@{}p{4cm}p{5cm}p{3cm}@{}}
\toprule
\textbf{Question} & \textbf{Why Important} & \textbf{Method Needed} \\
\midrule
Quality heterogeneity & Does quality of first experience matter? & RCT with quality variation \\
Frequency threshold & Where does learning $\to$ reinforcement? & Cross-domain comparison \\
Negative experiences & Do bad first experiences create aversion? & Natural experiments \\
Social learning & Can vicarious experience substitute? & Network experiments \\
Decay dynamics & Does learning effect decay over long horizons? & Long-term follow-ups \\
\bottomrule
\end{tabular}
\end{table}


%%%%%%%%%%%%%%%%%%%%%%%%%%%%%%%%%%%%%%%%%%%%%%%%%%%%%%%%%%%%%%%%%%%%%%%%%%%%%%%
\section{References}
\label{sec:lrn-references}
%%%%%%%%%%%%%%%%%%%%%%%%%%%%%%%%%%%%%%%%%%%%%%%%%%%%%%%%%%%%%%%%%%%%%%%%%%%%%%%

All references for this appendix are maintained in the master bibliography (\texttt{bcm\_master.bib}), Section ``APPENDIX HF: HABIT FORMATION \& EXPERIENCE GOODS''.

\textbf{Key Citations:}
\begin{itemize}[nosep]
    \item \texttt{Nelson1970}: Experience goods foundation
    \item \texttt{PAP-becker1988theory}: Rational addiction benchmark
    \item \texttt{PAP-ackerberg2003advertising}: Advertising and experience goods
    \item \texttt{PAP-dupas2014short}: Subsidies and learning
    \item \texttt{jones2025habits}: Key habit formation identification
    \item \texttt{PAP-banerjee2021challenges}: Preference discovery
    \item \texttt{hussam2022rational}: Frequency moderation
    \item \texttt{allcott2014short}: Persistence estimation
\end{itemize}


%%%%%%%%%%%%%%%%%%%%%%%%%%%%%%%%%%%%%%%%%%%%%%%%%%%%%%%%%%%%%%%%%%%%%%%%%%%%%%%
\section{Cross-References}
\label{sec:lrn-crossref}
%%%%%%%%%%%%%%%%%%%%%%%%%%%%%%%%%%%%%%%%%%%%%%%%%%%%%%%%%%%%%%%%%%%%%%%%%%%%%%%

\noindent\textit{Related Appendices:}
Appendix HF (PREDICT-HABIT),
Appendix UNMAPPED_PBB (FORMAL-PROBABILITY),
Appendix UNMAPPED_AU (CORE-AWARE),
Appendix UNMAPPED_BBB (CORE-WHERE),
Appendix XX (METHOD-EVIDENCE),
Appendix G (Glossary).

\noindent\textit{Related Chapters:}
Chapter 5 (Awareness),
Chapter 9 (Context),
Chapter 14 (Interventions),
Chapter 16 (Probability).

\noindent\textit{Master Bibliography:} \texttt{bcm\_master.bib}, Section ``APPENDIX HF: HABIT FORMATION \& EXPERIENCE GOODS''

\nocite{bcm_master}

% =============================================================================
% END OF APPENDIX XIX
% =============================================================================

\end{document}
